\insertImage{3c1da1ab-a826-4f5f-b681-6318fb79aca0.png}{Aquarelle d'une fenêtre de bureau moderne ouverte sur un environnement luxuriant et fleuri.}


\part*{Introduction}
\addcontentsline{toc}{part}{Introduction}

\textit{Libérer l'entreprise.} Au premier abord, l'idée semble aussi radicale qu'elle est séduisante. Et pourtant, il s'agit d'une pratique de plus en plus courante dans le monde des affaires d'aujourd'hui. Les entreprises libérées, ces entités audacieuses qui cherchent à réinventer la façon dont le travail est organisé et effectué, ont fait l'objet d'une attention considérable ces dernières années. Certains voient en elles le futur du travail\footnote{Fleming, P. (2019). The Worst Is Yet to Come: A Post-Capitalist Survival Guide. Repeater Books.}, tandis que d'autres y voient une tendance de gestion passagère\footnote{Gandini, A. (2019). Labour process theory and the gig economy. Human Relations, 72(6), 1039–1056. https://doi.org/10.1177/0018726718792867}. Quoi qu'il en soit, la réalité est souvent plus nuancée.

\section*{Présentation du concept d'entreprise libérée}

\subsection*{Définition et origines}

L'entreprise libérée est un modèle organisationnel qui vise à donner aux employés la liberté et la responsabilité de prendre leurs propres décisions, dans le but de favoriser l'innovation, l'efficacité et la satisfaction au travail. Ce concept, qui remonte à plusieurs décennies, a été popularisé par le professeur Isaac Getz de l'ESCP Europe dans son livre "Freedom, Inc."\footnote{Carney, B. M., \& Getz, I. (2009). Freedom, Inc.: Free Your Employees and Let Them Lead Your Business for Greater Success. Crown Business.}.

Il n'y a pas de définition unique et universellement acceptée de l'entreprise libérée, car chaque organisation est unique et adapte le concept à sa propre situation. Cependant, la plupart des entreprises libérées partagent certaines caractéristiques communes, notamment une décentralisation de l'autorité, une communication ouverte et transparente, une confiance et un respect mutuels, et un accent mis sur l'autonomie et l'initiative individuelle.

\subsection*{Principes et valeurs}

Les principes et valeurs qui sous-tendent le modèle de l'entreprise libérée sont enracinés dans des théories et des pratiques de gestion plus larges, y compris le management participatif, la théorie X et Y de Douglas McGregor, et les principes de responsabilité et d'autonomie prônés par Peter Drucker et d'autres gourous de la gestion\footnote{Drucker, P. (1954). The Practice of Management. Harper \& Brothers.}.

Dans une entreprise libérée, les employés sont non seulement autorisés, mais encouragés à prendre des initiatives et à faire preuve de créativité. La communication est ouverte et bidirectionnelle, avec un accent mis sur le partage des informations et des connaissances. La confiance est considérée comme un élément fondamental de la culture organisationnelle, et les managers sont vus comme des facilitateurs plutôt que comme des commandants.

\subsection*{Exemples d'entreprises libérées}

Des entreprises de toutes tailles et de tous secteurs ont adopté le modèle de l'entreprise libérée. Parmi les exemples les plus connus, on trouve FAVI, une entreprise française de sous-traitance automobile dirigée par Jean-François Zobrist\footnote{Zobrist, J.-F. (2012). La belle histoire de FAVI: L'entreprise qui croit que l'homme est bon. Humanisme \& Organisations.}, et Valve, un développeur américain de jeux vidéo célèbre pour sa structure organisationnelle plate et son manque de hiérarchie formelle\footnote{Laloux, F. (2014). Reinventing Organizations: A Guide to Creating Organizations Inspired by the Next Stage of Human Consciousness. Nelson Parker.}.

Ces entreprises ont été saluées pour leur capacité à stimuler l'innovation, à améliorer la satisfaction des employés et à obtenir des résultats financiers solides. Cependant, elles ont également dû faire face à des défis importants, notamment en matière de coordination, de contrôle et de gestion du changement.

Dans ce livre, nous explorerons ces défis et les leçons que nous pouvons en tirer. Nous examinerons également les critiques et les controverses qui entourent le modèle de l'entreprise libérée, et nous envisagerons des alternatives et des perspectives d'évolution.

En fin de compte, notre objectif est de fournir une vue d'ensemble équilibrée et approfondie de ce phénomène passionnant et complexe. Nous espérons que ce livre vous aidera à comprendre ce que signifie vraiment "libérer l'entreprise" - et peut-être à déterminer si c'est une voie que vous, ou votre organisation, souhaitez emprunter.
